Em $K=1$, 24 das 28 instâncias admitem atendimento completo. A comparação de distância média entre métodos resulta em 27 vitórias da tabu, 0 do VND e 1 empate. O fallback foi usado em 0 das 1680 construções desse caso.

Em $K=2$, 28 das 28 instâncias admitem atendimento completo. A comparação de distância média entre métodos resulta em 28 vitórias da tabu, 0 do VND e 0 empates. O fallback foi usado em 0 das 1680 construções desse caso.

As quatro inviabilidades em $K=1$ são I1 G2, I2 G2, I2 G3 e I2 G5; essa classificação depende de $K$. Em I2 G2 (109 candidatos, seis tipos e cinco escolas), o atendimento certificado passa de 96 em $K=1$ para 109 em $K=2$. A distância parcial de $K=1$ não mede a mesma tarefa da distância completa de $K=2$; portanto, não calculamos melhoria percentual entre esses dois casos.

O VND foi executado até o ótimo local em todas as execuções: o teto de segurança de 1000000 iterações foi atingido 0 vezes; a trajetória média teve 863 iterações e a mais longa 7016. A tabu atingiu o teto de 1000 iterações em 776 das 1680 execuções; nas demais, parou após 200 iterações sem melhora. A melhoria média sobre a construção foi de 12,87\% no VND e 22,22\% na tabu, com tempo médio de busca de 0,519 s e 1,628 s por execução, respectivamente.

O pior desempenho relativo do VND ocorreu em I2 G8 com $K=2$: a distância média final (329,64 km, desvio padrão 112,58) ficou 209,60\% acima da melhor solução conhecida (106,47 km, mínimo entre as 60 execuções dos dois métodos), e a pior semente ficou 513,41\% acima. A tabu obteve média de 107,87 km na mesma instância.

Todas as 6720 soluções iniciais/finais respeitaram capacidade e compatibilidade e atingiram o atendimento certificado. As contagens de vitórias são descritivas e comparam os métodos somente dentro da mesma instância, do mesmo K e do mesmo atendimento; não estabelecem significância estatística.

A pré-checagem é executada uma vez por lote de 30 sementes, separadamente para cada método e K. A tabela de tempos inclui esse custo integral em P e o amortiza em Total: $T_{\mathrm{total}}=(T_{\mathrm{leitura}}+P)/30+T_{\text{construção}}+T_{\mathrm{busca}}$. O custo do lote é 30 vezes esse total médio; não se trata de medição isolada com CPU exclusiva.

Ambiente: macOS (arm64), Python 3.9.6, C++17 com Apple Clang e otimização \texttt{-O2}, NumPy 1.24.2 e SciPy 1.11.3. VND e tabu foram executados simultaneamente, com quatro processos cada (oito processos em oito núcleos); os tempos incluem competição por CPU. Os manifestos registram entradas, binário e parâmetros. O prefixo do hash do binário é \texttt{a4274b9a6a7ffae1}.
